%% §9 Conclusion — D-05 practitioner decision tree + winner_per_regime figure
%% + V2-01..V2-06 limitations/future work.
\section{Conclusion}\label{sec:conclusion}

We compared six cache eviction policies (LRU, FIFO, CLOCK, S3-FIFO, SIEVE,
W-TinyLFU) on two real public-records API workloads --- Congress.gov and
CourtListener --- using a 5-seed replay-Zipf model with Welch's $t$-test
significance testing. W-TinyLFU wins 11 of 12 miss-ratio cells overall, but
the headline finding is more nuanced than that count suggests: on Congress's
near-uniform raw trace, SIEVE's visited-bit is functionally equivalent to
W-TinyLFU's Count-Min Sketch across the entire $\alpha$ range from $0.6$ to
$1.5$ (\S\ref{sec:mechanism}). The admission-filter benefit is
workload-gated, not policy-inherent.

\subsection{Winner per regime}

\begin{figure}[h]
\centering
\includegraphics[width=\textwidth]{../results/compare/figures/winner_per_regime_bar.pdf}
\caption{Winner-per-regime summary. Each bar shows the winning policy's
miss-ratio (lower is better) in that regime on each workload. Base policies
only; ablation variants excluded per \S\ref{sec:ablations}.}
\label{fig:winner_bar}
\end{figure}

\begin{table}[h]
\centering
\footnotesize
\renewcommand{\arraystretch}{0.92}
\begin{tabularx}{\textwidth}{@{}lXXXX@{}}
\toprule
\textbf{Regime} & \textbf{Congress Winner} & \textbf{Cong Miss} & \textbf{Court Winner} & \textbf{Court Miss} \\
\midrule
Small Cache (\texttt{cache\_frac=0.001})            & W-TinyLFU & 0.869 & W-TinyLFU & 0.831 \\
High Skew ($\alpha \in \{1.0, 1.1, 1.2\}$)          & SIEVE     & 0.351 & W-TinyLFU & 0.386 \\
Mixed Sizes (Court byte-MRC, single-seed)           & N/A       & ---   & W-TinyLFU & 0.284 \\
OHW Regime (\texttt{cache\_frac=0.01})              & W-TinyLFU & 0.712 & W-TinyLFU & 0.728 \\
\bottomrule
\end{tabularx}
\caption{Winner per regime across both workloads. Single-seed for the Mixed
Sizes Court byte-miss-ratio row (multi-seed byte-MRC aggregation deferred to
v2; see V2-05 below). Mixed Sizes on Congress is N/A because Congress has
near-uniform object sizes (Table~\ref{tab:workload}).}
\label{tab:winner_regime}
\end{table}

\subsection{A practitioner decision tree}

The most actionable takeaway from this study is a small flowchart for picking
an eviction policy for a public-records API workload. Use this decision tree
given a raw trace's characteristics:

\begin{itemize}\setlength{\itemsep}{2pt}
  \item \textbf{If your raw trace has $\alpha_{\text{raw}} < 0.5$ and a
  near-uniform size distribution (Congress-like):} SIEVE or W-TinyLFU are
  statistically tied across $\alpha \in [0.6, 1.5]$
  (\S\ref{sec:mechanism}, Table~\ref{tab:mechanism_gap}). Prefer SIEVE for
  its simplicity --- it has no admission filter, no frequency sketch, no
  Count-Min Sketch, just a visited-bit and a hand pointer. Implementation
  and operational-debugging cost is lower for identical miss-ratio
  performance.
  \item \textbf{If your raw trace has $\alpha_{\text{raw}} > 0.8$ and a
  long-tail size distribution (Court-like):} \textbf{W-TinyLFU by $\sim 5$pp
  at every $\alpha$.} The Count-Min Sketch's frequency gradient extracts
  information SIEVE's binary visited-bit cannot. The implementation cost is
  real but modest (a CMS, a Doorkeeper Bloom filter, an admission-filter
  comparison); the performance gap is stable enough to justify it.
  \item \textbf{If you have a catastrophic size outlier (Court-like):}
  W-TinyLFU's admission filter matters most at small cache sizes, where a
  single one-hit-wonder admission can evict 5--10\% of capacity. Report
  \emph{both} object and byte miss-ratios with 5-seed CI bands; a
  single-seed byte-miss number is dominated by which Zipf rank the outlier
  lands at (\S\ref{sec:byte-mrc}).
  \item \textbf{Small caches amplify everything:} if \texttt{cache\_frac}
  $< 0.01$, admission quality dominates (W-TinyLFU beats LRU by 9--13pp on
  both workloads in the Small Cache regime). If \texttt{cache\_frac} $>
  0.1$, all reasonable policies converge; eviction algorithm choice stops
  mattering.
\end{itemize}

\subsection{Limitations and future work}

Six items are acknowledged but deferred to a future revision. None of them
changes any finding reported here; together they extend the rigor envelope.

\begin{itemize}\setlength{\itemsep}{2pt}
  \item \textbf{Caffeine trace cross-validation (V2-01):} running our
  W-TinyLFU implementation on a Caffeine-published benchmark trace would
  provide a $\pm 2\%$ agreement check against the reference Java
  implementation \cite{caffeine2024}; deferred due to benchmark-trace access
  friction.
  \item \textbf{LHD or AdaptSize as a 7th policy (V2-02):} size-aware
  eviction would likely behave quite differently on Court's long-tail size
  distribution and would provide a direct counterfactual for the
  \S\ref{sec:byte-mrc} outlier-variance story.
  \item \textbf{Third trace domain (V2-03):} a SEC EDGAR collector would
  complete the public-records trilogy and test whether the Congress-like /
  Court-like dichotomy generalizes or fragments across more diverse
  workload traits.
  \item \textbf{Multi-seed $\alpha \in \{1.3, 1.4, 1.5\}$ on Congress
  (V2-04):} the single-seed crossover data in \S\ref{sec:mechanism} is
  suggestive but not CI-backed. Cheap to add ($\sim 30$\,s wall-clock per
  seed $\times$ 3 new $\alpha$ values $\times$ 5 seeds).
  \item \textbf{Multi-seed byte-MRC aggregation (V2-05):} the per-seed CSVs
  at \texttt{results/compare/multiseed/court/} already record
  \texttt{byte\_miss\_ratio}; the Phase~5 aggregator only consumes
  \texttt{miss\_ratio}. A one-line change to
  \texttt{compare\_workloads.py::aggregate} would retire the
  ``single-seed'' hedge on the Mixed Sizes row of
  Table~\ref{tab:winner_regime}.
  \item \textbf{Code review WR-01 / WR-02 (V2-06):} two latent-coupling
  warnings in \texttt{plot\_results.py} (BASE\_POLICIES filter missing on
  winner-bar; regime dispatch missing else clause). Neither manifests in
  the current pipeline; both would bite future work that adds ablation
  variants to the multi-seed sweep.
\end{itemize}

These items do not change any finding reported here; they extend the rigor
envelope for future work.
